We have presented \Gweek{}, an implementation of a functional logic programming
language whose design is derived from the algebraic theory of FLP effects. The
key insight driving the design is that the characteristic constructs of
FLP---nondeterministic choice, logical variables, and equational
constraints---are algebraic effects in the sense of \textcite{plotkin_2003}, and
can be given a clean account in the CBPV framework of \textcite{levy_2003}.

The implementation demonstrates that the algebraic theory is not merely a
theoretical curiosity: it provides direct guidance for the design of the
abstract machine, the equational laws used by the optimiser, and the
understanding of when and how effects interact. The suspension mechanism, in
particular, arises naturally from the CBPV treatment of sequencing, and its
consequences for programming---the need to force constraints, the interaction
between laziness and nondeterminism, the dramatic effect of interleaving
generation and testing---are illuminated by the equational theory.

Several directions for future work present themselves. First, the language could
be extended with \emph{encapsulated search}: a primitive that collects all
results of a nondeterministic computation into a value. As discussed in
\parencite{jones_2026}, this requires moving beyond algebraic effects to
\emph{effect handlers}, and raises interesting questions about the interaction of
encapsulation with the commutativity and idempotence of choice. Second, the type
system could be extended with polymorphism and type classes, enabling more
expressive programs. Third, the implementation could benefit from more
sophisticated search strategies, such as constraint propagation or learned
heuristics, which would bring \Gweek{} closer to the state of the art in
constraint programming. Finally, the strict evaluation mode suggests an
interesting spectrum between lazy and strict FLP semantics, whose theoretical
and practical implications deserve further study.
